\section{A quadratic model with its Frobenius action}
\label{mod:section}

We now replace each Frobenius-fixed retained component by a derived
quadratic moment-map model, retaining its evaluation bundles and the map
defining nilpotent singular support. The main issue is to extend
coordinates constructed at its closed semisimple orbit over the formal
invariant base.

Fix an abstract field isomorphism $\E\simeq\C$. All absolute values in
this section refer to this isomorphism. An invertible linear operator has
\emph{absolute weight $w$} if all its eigenvalues have modulus $q^{w/2}$.
This terminology does not assert that the operator is semisimple. It is
distinct from both cohomological degree and the weights of an
$\mathfrak{sl}_2$-triple.

The spectral stack remains the prime union $Z$ from
Section~\ref{inp:section}. Its components form a categorical coproduct.
In the twisted bar construction, an arrow string closes only in a
component preserved by $\Phi$. Consequently
\begin{equation}
\label{mod:component-trace}
 \Tr(\Phi,\IndCoh_Y(Z))
 \simeq\bigoplus_{Z_\sigma\text{ preserved by }\Phi}
       \Tr(\Phi,\IndCoh_Y(Z_\sigma)).
\end{equation}
The sum is a direct sum, including when there are infinitely many
components. There are no arrows between distinct components; this also
proves their orthogonality for the Hom pairings used below.

Choose such a component and its unique closed semisimple parameter
$\sigma$. Choose a pure Weil lift as in Theorem~\ref{inp:pure-lift},
and put
\[
 J=\Aut(\sigma),\qquad \mathfrak j=\Lie(J),\qquad
 V=H^1(X_{\overline{\mathbb F}_q},\ghat_\sigma).
\]
The group $J$ is reductive; it need not be connected. Curve duality and
an invariant form identify $H^2(\ghat_\sigma)$ with
$\mathfrak j^*(-1)$ and give a $J$-invariant symplectic form $\omega$ on
$V$, after forgetting its Tate factor. The moment map is
\begin{equation}
\label{mod:moment}
 \mu:V\longrightarrow\mathfrak j^*,\qquad
 \langle\mu(x),z\rangle=\tfrac12\omega(zx,x).
\end{equation}
Its zero fiber will always mean its \emph{derived} zero fiber. In
particular no regular-sequence assumption on the components of $\mu$
is being made.

\subsection{What is recovered from completion at the closed orbit}

We first explain why constructing a model near $\sigma$ suffices here.
The invariant base must be retained in this argument. Choose a closed
point $v$ of degree $d$ and frames at its $d$ geometric points. Write
$K=\Ghat^d$. The framed component supplied by
Theorem~\ref{inp:components} is an ind-affine derived scheme over a
formal invariant base with reduced scheme a point. Its unique closed
$K$-orbit is $K/J$. It has affine almost-finite-type stages over
nilpotent thickenings of that invariant point. At each such stage the
classical coordinate ring is noetherian, its invariant ring is
Artinian, and its cohomology sheaves in any bounded Postnikov range are
coherent. Adding the other frames to a one-frame presentation is an
affine operation. These are precisely the finite-stage properties we
use.

\begin{lemma}[Finite covariants and completion]
\label{mod:finite-covariants}
Let $K$ be a reductive complex algebraic group, possibly disconnected,
and let $A$ be a finitely generated $K$-algebra. Suppose $A^K$ is
Artinian and $\Spec A$ has a unique closed $K$-orbit $O$, with ideal
$I$. For an equivariant finite $A$-module $M$ and a finite-dimensional
$K$-representation $D$, the inverse system
\[
 \Hom_K(D,M/I^nM)
\]
is eventually constant with value $\Hom_K(D,M)$. In particular the
rational, locally $K$-finite part of the $I$-adic completion recovers
$M$, and recovers $A$ together with its multiplication when $M=A$.
\end{lemma}

\begin{proof}
Choose a finite-dimensional $K$-stable generating space $U\subset M$.
The surjection $A\otimes U\to M$ induces a surjection on
$\Hom_K(D,-)$, by reductivity. Finite generation of the covariants
of $A$ over $A^K$ \cite[Lemma~2.3]{Brion} therefore makes
$M_D=\Hom_K(D,M)$ finite over $A^K$, hence finite-dimensional.
This applies to disconnected $K$: its identity component is reductive
and its component group is finite. Let $N=\bigcap_{n\geq1}I^nM$. It is a finite
equivariant submodule because $A$ is noetherian. At every point of $O$
its localization vanishes by Krull intersection \cite[Tag~00IP]{Stacks}. If $N$ were nonzero,
its nonempty invariant closed support would contain a closed orbit,
hence would meet $O$. Therefore $N=0$.

The descending subspaces $\Hom_K(D,I^nM)$ of $M_D$ have zero
intersection and thus are eventually zero. Reductivity in
characteristic zero makes $\Hom_K(D,-)$ exact, so
\[
 \Hom_K(D,M/I^nM)
 =M_D/\Hom_K(D,I^nM).
\]
This proves the assertion. Every rational representation is the union
of its finite-dimensional subrepresentations; applying the assertion
to all $D$ therefore recovers $M$. Applied also to finite tensor
products of representation spaces, it recovers multiplication and
every equivariant finite-presentation map.
\end{proof}

For example, let $\Gm$ act with weights $1,-1$ on the coordinate
functions $x,y$, and put
\[
 A_N=\C[x,y]/((xy)^N),\qquad
 A_N^{\Gm}=\C[t]/(t^N),\quad t=xy,\quad N\ge1.
\]
The origin is the unique closed orbit. Each weight space of $A_N$ is
finite-dimensional, so the locally $\Gm$-finite part of its
$(x,y)$-adic completion recovers $A_N$. The ordinary inverse limit
also contains arbitrary formal series such as $\sum_{m\ge0}x^m$,
which are not locally finite. Thus rational recovery retains the
nonzero points on both axes, not only the closed orbit. As $N$
varies, these stages give the formal neighborhood of $xy=0$,
including the formal invariant parameter $t$. The next lemma carries
out this recovery for derived algebras and their maps.

\begin{lemma}[Restoring the invariant base]
\label{mod:reconstruction}
For the framed component above, an equivariant isomorphism of its
completion along $K/J$ with another such completed model extends
uniquely over the formal invariant bases to an isomorphism of their
invariant-fiber completions. The same assertion holds for equivariant
maps and their full mapping spaces, and for finite equivariant
presentation data. Here the extension is unique up to a contractible
space of choices.
\end{lemma}

\begin{proof}
We give the derived argument, since recovering coordinate rings alone
would not recover homotopies. The proof has three stages: finite
covariants recover each affine stage and its mapping spaces; trivial
covariants recover the invariant base; derived Koszul thickenings then
give compatible extensions over that base.
Write $\mathcal R_K$ for the derived
category of rational $K$-representations. All limits in this paragraph
are taken in $\mathcal R_K$, not in the category of vector spaces with
an arbitrary, possibly nonrational, $K$-action.

\emph{Recovery at a fixed invariant-base stage.}
Let $A$ be its connective coordinate algebra, viewed as a commutative
algebra in $\mathcal R_K$. The precise hypotheses are that $H^0(A)$
is finitely generated, $H^0(A)^K$ is Artinian,
$\Spec H^0(A)$ has a unique closed orbit $O$, and every $H^i(A)$
is finite over $H^0(A)$. Choose a
finite-dimensional $K$-stable generating space for the ideal of $O$
in $H^0(A)$, lift its generators to degree-zero cycles, and extend
them to a finite-dimensional stable generating space for $H^0(A)$.
Reductivity allows these lifts to be equivariant. This gives a map
from a polynomial $K$-algebra $P$ to $A$, surjective on $H^0$.
Let $I_P$ be the ideal generated by the chosen orbit equations and put
\[
 A_n=A\otimes_P^{\mathbf L}P/I_P^n.
\]
These derived orbit jets compute completion along $O$; the
module-level comparison is \cite[Tag~0BKH]{Stacks}.
The resulting formal neighborhood is independent of the generating
presentation \cite[Theorem~A.1.3(c) and \S A.1.5]{AGKRRV1}.

Here is the cohomological calculation at these jets. Every $H^i(A)$
is a finite $P$-module. Noetherian Artin--Rees makes the positive-Tor
systems
$\{\operatorname{Tor}^P_j(H^i(A),P/I_P^n)\}_n$, $j>0$,
pro-zero. Since $P$ has finite global dimension, the hyper-Tor
calculation has finite width in every cohomological degree, even if
$A$ is unbounded below. Consequently
\[
 \{H^i(A_n)\}_n\simeq
 \{H^i(A)/I_P^nH^i(A)\}_n
\]
as pro-modules. This is the coherent derived-completion calculation
of \cite[Tags~00MA, 0A06, and~0BKH]{Stacks}, applied here over the
ordinary noetherian polynomial ring $P$.
For every finite-dimensional $K$-representation $D$,
Lemma~\ref{mod:finite-covariants} and exactness of covariants therefore
give a natural pro-isomorphism
\begin{equation}
\label{mod:derived-covariants}
 \{H^i\RHom_K(D,A_n)\}_n
 \simeq H^i\RHom_K(D,A),
\end{equation}
with a constant system on the right. In particular these systems
have zero $\lim^1$. The functors $\RHom_K(D,-)$ commute with limits
and jointly detect equivalences in $\mathcal R_K$. Applying the
Milnor exact sequence to~\eqref{mod:derived-covariants} proves
\begin{equation}
\label{mod:rational-derived-completion}
 A\xrightarrow{\sim}\operatorname*{holim}_n A_n
       \quad\text{in }\operatorname{CAlg}(\mathcal R_K).
\end{equation}
Indeed the forgetful functor from commutative algebras creates limits
and detects equivalences. Thus this is a natural equivalence of
algebras, not merely an identification of their cohomology. It gives,
for every rational equivariant coordinate algebra $B$, an equivalence
of full mapping spaces
\begin{equation}
\label{mod:mapping-spaces}
 \operatorname{Map}_{\operatorname{CAlg}(\mathcal R_K)}(B,A)
 \simeq
 \operatorname*{holim}_n
 \operatorname{Map}_{\operatorname{CAlg}(\mathcal R_K)}(B,A_n),
\end{equation}
compatible with composition. No index $n$ uniform in $i$ or $D$ has
been asserted or used. The rational limit is essential: as in the preceding example, the
ordinary vector-space limit would retain arbitrary formal series.
The same calculation applies to an $A$-module with coherent
cohomology. Finite equivariant module presentations and their maps
are recovered by testing their generating representation spaces;
their relations and coherent homotopies are recovered by the
corresponding derived mapping spaces.

\emph{Recovery of the invariant base.}
The formal coarse base is constructed from the system of trivial
covariants of the framed affine stages
\cite[\S\S5.3.5--5.3.6]{AGKRRV1}. We do not identify an arbitrarily
chosen base stage with the spectrum of its pullback's invariant
algebra. The trivial case of~\eqref{mod:derived-covariants} instead
recovers this entire system. On classical invariant functions the
topologies are cofinal: if $\mathfrak m$ is the invariant ideal of
definition, then $\mathfrak m\subset I$, and, after passage to a
fixed nilpotent invariant-base stage,
Lemma~\ref{mod:finite-covariants} makes the trivial covariants of high
powers of $I$ vanish. Applying this in every bounded range gives the
corresponding comparison for the derived invariant systems.
An equivariant map of orbit completions therefore gives the formal
base map as well.

\emph{Factorization and extension.}
The noetherian formal-affine presentation of that base uses finitely
many invariant generators $t_1,\ldots,t_r$ of an ideal of definition.
The noetherian and coherent presentation is supplied by
\cite[Theorem~A.1.3(b),(c) and Remark~A.1.4]{AGKRRV1}.
Its infinitesimal stages may be taken to be the derived Koszul
thickenings
\begin{equation}
\label{mod:koszul-base-stages}
 R_N=R/\!/(t_1^N,\ldots,t_r^N)
   :=R\otimes_{\C[u_1,\ldots,u_r]}^{\mathbf L}\C,
 \qquad u_j\longmapsto t_j^N.
\end{equation}
These are precisely the formal-affine stages used in
\cite[\S7.1 and \S\S7.9.8--7.9.10]{AGKRRV1}. The relevant
framed base changes are affine and almost of finite type.

Restrict a map of orbit completions to the completed orbit of one
source stage $A$. Compose it with the target's map to its formal base
and then with the structural map to $\Spec R$.
Equation~\eqref{mod:mapping-spaces}, applied to the base algebra $R$
with trivial $K$-action, recovers a map $R\to A$ and its homotopies.
The images of the $t_j$ vanish at the reduced invariant point, so
they are nilpotent in the Artinian algebra $H^0(A)^K$. Choose a common power $N$ and invariant nullhomotopies of
these $t_j^N$. Such nullhomotopies exist because vanishing in $H^0$
means being a boundary; they may be chosen invariant by reductivity.
The universal property of~\eqref{mod:koszul-base-stages} now factors
this base map through a derived Koszul stage. Pulling back the target
framed component along that stage factors the completed-space map
through an affine target stage. Nilpotence
alone would not justify factorization through an ordinary quotient.
If $B_N$ is the coordinate algebra of the resulting affine target
stage, the completed-stage map is a compatible point of the
right-hand side of~\eqref{mod:mapping-spaces}, with $B=B_N$.
It therefore extends to a map on the source stage $A$, with all its
homotopies and composition data.

For clarity, the choices of powers and nullhomotopies do not add
extra maps. The Koszul stages are cofinal tests for a formal
neighborhood. One can check this directly on mapping spaces: above
a fixed map with nilpotent $t_j$, the choices of nullhomotopy are
torsors for the corresponding additive loop spaces, and increasing
$N$ to $M$ acts on their differences by multiplication by
$t_j^{M-N}$. Nilpotence kills every such homotopy group after passing
far enough in this filtered system. Its nonempty space of choices
is therefore contractible. This verifies cofinality for maps and
homotopies, without interchanging a filtered colimit with a
Postnikov limit. Equation~\eqref{mod:mapping-spaces} makes the
extended maps on different source stages compatible. Applying it
also to inverse maps and to their compositions completes the proof.
\end{proof}

Thus completion along the closed orbit is used to \emph{construct}
coordinates. It is not substituted for completion along the invariant
null fiber. For example, nonclosed points of that null fiber are
retained by Lemma~\ref{mod:reconstruction}.

\subsection{Simultaneous linearization}

Here is the elementary nonresonance calculation that removes the
higher relations in the local deformation problem.
The classical Frobenius-compatible quadratic hull in the smooth proper,
pure weight-zero setting is established in
\cite[Theorem~2.11 and Corollary~2.12]{Pridham}; its proof also uses
Jordan decomposition on finite jets. We retain the derived presentation
and equivariant data, then apply Lemma~\ref{mod:reconstruction} to
restore the formal invariant base.

\begin{lemma}[Nonresonance with Jordan blocks]
\label{mod:nonresonance}
Let a minimal complete free graded-commutative algebra have generators
$V^*$ in cohomological degree $0$ and $R^*$ in degree $-1$. Suppose a
continuous graded automorphism has linear parts of absolute weights
$-1$ and $-2$ on these two generator spaces. It is conjugate, by a
formal graded coordinate change tangent to the identity, to its
linear part. If it commutes with a minimal differential, that
differential, in the resulting coordinates, sends $R^*$ to quadratic
polynomials in $V^*$. These assertions hold equivariantly for a
reductive group normalized by the automorphism.
\end{lemma}

\begin{proof}
Filter coordinate changes by total generator order. At order $n\geq2$
the error in a degree-zero generator belongs to
$\Hom(V^*,\Sym^nV^*)$. All eigenvalues of the conjugacy operator
``linear target action times inverse linear source action'' on this
space have modulus $q^{-(n-1)/2}$, hence differ from $1$. An error in
a degree-minus-one generator contains exactly one degree-minus-one
factor and $n-1$ degree-zero factors. The corresponding operator on
\[
 \Hom(R^*,R^*\otimes\Sym^{n-1}V^*)
\]
has the same modulus $q^{-(n-1)/2}$. In either case its difference
from the identity is invertible. This assertion follows from its
eigenvalues and holds on generalized eigenspaces; it requires no
diagonalization. Solve the conjugacy equation at this order and
continue. The corrections converge in the generator-adic topology.

In the equivariant case the errors and corrections lie in the
equivariant subspaces of these finite-dimensional jet spaces. The
conjugacy operator preserves these subspaces and its inverse does
also, so the same construction is equivariant. Transport the
differential along the coordinate changes. The differential of a
degree-zero generator is zero by degree. That of a degree-minus-one
generator has no constant or linear term by minimality. Since the
differential commutes with the now linear automorphism, an order-$n$
term can occur only if its weight $-n$ agrees with $-2$. Thus only
$n=2$ remains.
\end{proof}

\begin{theorem}[The framed quadratic model]
\label{mod:model}
\label{mod:quadratic-model}
The retained component $Z_\sigma$, with its Frobenius action,
evaluation bundles at the chosen cyclic tuple, and singular-support
map, admits the model
\begin{equation}
\label{mod:quadratic-equation}
 Z_\sigma\simeq
 \bigl[(\mu^{-1}(0))^{\wedge}_{\mathrm{unst}}/J\bigr].
\end{equation}
Here $\mathrm{unst}$ is the inverse image of the origin under
$V\to V/\!/J$, and completion is along this invariant null fiber.
In this model:
\begin{enumerate}[label=\textup{(\roman*)}]
\item the singularity scheme has points $(x,z)$ satisfying
      $\mu(x)=0$ and $zx=0$, with $z\in\mathfrak j$;
\item the support condition defined by $Y$ is the condition that the
      image of $z$ under $\mathfrak j\hookrightarrow\ghat$ belongs
      to $Y$;
\item every evaluation representation of $K=\Ghat^d$ is the bundle
      associated to its restriction along $J\hookrightarrow K$;
      Frobenius has constant cyclic transport on these bundles;
\item after removing the scalars $q^{1/2}$ on deformation points and
      $q^{-1}$ on singular points, Frobenius is a linear
      automorphism of the moment-map data, normalizing $J$ and the
      framed representation labels. Its eigenvalues on these
      normalized linear data have modulus one.
\end{enumerate}
The meaning of Frobenius on Hom complexes in \textup{(iv)} is fixed
explicitly in Lemma~\ref{mod:hom-action} below.
\end{theorem}

\begin{proof}
We first work in the completion of the framed component along $K/J$.
The smooth affine orbit admits an equivariant formal retraction:
lift the identity retraction through successive nilpotent
neighborhoods, using smoothness, and average the lifting torsors
using reductivity. The fiber is a formal slice with stabilizer $J$;
the completed framed space is its induction from $J$ to $K$.
Equivalently, this follows by splitting the orbit tangent directions
and comparing tangent complexes at every lifting step.

The transverse tangent and obstruction spaces are respectively
$V=H^1(\ghat_\sigma)$ and $H^2(\ghat_\sigma)$. Quasi-smoothness
therefore gives a minimal presentation by complete free coordinates
dual to these spaces in degrees $0$ and $-1$. One obtains this
presentation from a free resolution by removing the contractible
linear pairs; the splittings can be chosen $J$-equivariantly.

We recall the second-order obstruction calculation to fix its coefficient
and sign, working over $E$ before the chosen coefficient transport.
Use the simplicial descent model for local systems underlying
\eqref{inp:tangent}, computed by hypercover cochains with the
$\ell$-adic limit retained. Its adjoint cochains compute
$R\Gamma(\overline X,\ghat_\sigma)$, and its cup bracket combines the
cup product with the Lie bracket of the adjoint coefficients. For a
formal parameter $t$, write a perturbation of the descent data modulo
$t^3$ as $g_{ij}\exp(t x_{ij}+t^2 y_{ij})$, with the cochains twisted by
the original transition data $g_{ij}$. Expanding the cocycle equation
by the Baker--Campbell--Hausdorff formula gives
\[
 dx=0,\qquad dy=-\tfrac12[x,x].
\]
Thus the quadratic obstruction is the class of $\tfrac12[x,x]$ in
$H^2(\overline X,\ghat_\sigma)$. This calculation uses a cochain model
computing the deformation problem, rather than an assumed finite
trivializing cover. Invariance of the coefficient pairing and graded
commutativity give, with the duality and sign convention of
\eqref{mod:moment},
\[
 \bigl\langle\tfrac12[x,x],z\bigr\rangle
       =\tfrac12\omega(zx,x)=\langle\mu(x),z\rangle.
\]
This identifies the quadratic obstruction with $\mu(x)$.

We next arrange that both the retraction and these coordinates respect
the chosen Frobenius transporter. Purity gives absolute point
weights $0,1,2$ on the orbit tangent, transverse tangent, and
obstruction spaces. A positive-transverse-order error in the
retraction has positive source weight and orbit target weight zero.
Its Frobenius conjugacy operator therefore has no eigenvalue $1$.
At every finite jet order it can be removed uniquely by solving the
same linear equation used in Lemma~\ref{mod:nonresonance}. The errors
are equivariant sections of the tangent bundle of the smooth orbit;
they are determined by their $J$-equivariant values at its base point.
The stabilizer transport and the cyclic permutation of frames have
weight zero. This constructs a Frobenius-compatible retraction and
an invariant slice.

Represent its single cyclic automorphism on the minimal free
presentation. Such a representative can be chosen equivariantly by
free lifting and reductivity. Its invertible linear part makes it an
automorphism of the complete graded algebra. There is no additional
relation to impose on the generator of an infinite cyclic action.
On formal coordinate functions use inverse transport; its generator
weights are $-1,-2$. Lemma~\ref{mod:nonresonance} makes this
automorphism linear and forces the transported differential to be
quadratic. A coordinate change tangent to the identity preserves
that quadratic part, so the differential is precisely $\mu$.
This proves the moment-map presentation on the orbit completion.
Lemma~\ref{mod:reconstruction} extends it over the invariant base,
giving~\eqref{mod:quadratic-equation}.

For its singularities one computes directly that
\[
 \langle d\mu_x(v),z\rangle=\omega(zx,v).
\]
Nondegeneracy of $\omega$ identifies
$\ker(d\mu_x^*)$ with $\{z:zx=0\}$, proving \textup{(i)}.
The singular evaluation map is fiberwise linear in $z$. At $x=0$
curve duality identifies it with the centralizer inclusion
$\mathfrak j\hookrightarrow\ghat$, with its duality Tate factor.
The target singular points have absolute weight $-2$. A correction
of positive degree $n$ in $x$, linear in $z$, would have weight
$n-2$, and cannot intertwine this Frobenius action. Comparison on
each finite jet, again valid for Jordan blocks, excludes every such
term. This proves \textup{(ii)}. Changing the invariant form only
rescales the simple factors of $\ghat$, which preserves every
nilpotent orbit.

Finally, a $K$-frame is exactly the datum trivializing all the
evaluation bundles in question. Write $\mathscr S$ for the formal
slice. The equivariant retraction identifies the framed completion
with $K\times^J\mathscr S$, with its map to $K/J$ given by
$[k,s]\mapsto kJ$. For a representation $D$ of $K$, quotienting the
trivial framed bundle by $K$ therefore gives the bundle on
$[\mathscr S/J]$ associated to $D|_J$. The retraction is
Frobenius-compatible, so the transport on its $K/J$ coordinate is
the closed-orbit transport, independent of $s$. The same is therefore
true of the cyclic transport on every evaluation bundle.
Purity on the chosen Weil
parameter gives modulus-one eigenvalues on every normalized frame
representation. A cyclic action can be tested on its $d$th power.
Together with purity on $H^1$ and $H^2$, this proves
\textup{(iii)} and \textup{(iv)}.
\end{proof}

\subsection{Which action contracts Hom complexes}

The next convention is essential: the category still has the
endomorphism $\Phi$ specified in the problem.

\begin{lemma}[Inverse transport on arrows]
\label{mod:frobenius}
\label{mod:hom-action}
Choose structures $\alpha_M:\Phi M\xrightarrow{\sim}M$ on compact
objects. On an arrow $h:M\to N$ put
\begin{align}
 T(h)&=\alpha_N\,\Phi(h)\,\alpha_M^{-1},\label{mod:T}\\
 S(h)&=\Phi^{-1}(\alpha_N^{-1}h\alpha_M)=T^{-1}(h).
 \label{mod:S}
\end{align}
The action in the bilinear Hom pairing is $S$. On the quadratic
model it gives absolute weight $-1$ to deformation functions and
absolute weight $+2$ to degree-two obstruction operators. Thus its
associated point weights on $(x,z)$ are $(1,-2)$, and
$W(x,z)=\langle z,\mu(x)\rangle$ has weight zero.

The graph identification in~\eqref{tr:graph-identification} turns
the original $\Phi$-graph bimodule into the graph of the transported
arrow action. In those coordinates the induced simultaneous bar
actions $T$ and $S$ are inverse. If the paracyclic homotopy is
$bH+Hb=1-T$, then
\begin{equation}
\label{mod:inverse-homotopy}
 b(-T^{-1}H)+(-T^{-1}H)b=1-S.
\end{equation}
In particular using contracting $S$-weights does not replace
$\Tr(\Phi,-)$ by $\Tr(\Phi^{-1},-)$.
\end{lemma}

\begin{proof}
The Frobenius compatibility in Section~\ref{inp:section} identifies
$\Phi$ with $f^!$, where $f$ is pullback of parameters by geometric
Frobenius of the curve. Since $f$ is an automorphism, this is its
usual pullback functor. A covariant trace structure is
$b:N\to\Phi N$, whereas a structure on the opposite input is
$a:\Phi M\to M$. Their pairing takes an arrow to
\[
 h\longmapsto\Phi^{-1}(b\,h\,a).
\]
With $b=\alpha_N^{-1}$ and $a=\alpha_M$ this is~\eqref{mod:S}.
It is the inverse of~\eqref{mod:T}, directly by composition.

For the signs, consider a coordinate and a relation generator with
$f^\#x=a_1x$ and $f^\#\epsilon=a_2\epsilon$, where
$|x|=0$, $|\epsilon|=-1$. Inverse transport on the endomorphisms of
the free module sends multiplication by $x$ to multiplication by
$a_1^{-1}x$. In a Koszul resolution the degree-two operator dual to
$\epsilon$ transforms by $a_2$ under the same transport. Thus
purity, $|a_1|=q^{1/2}$ and $|a_2|=q$, gives the stated signs.
The calculation is linear-algebraic and applies to generalized
eigenspaces. The polynomial operator dual to $\epsilon$ is a linear
function on the singular variable $z$; point weights are the
contragredient weights, giving $-2$ on $z$.

There is a distinction between ordinary arrows and raw graph
coefficients. In the left-graph convention for classes $M\to\Phi M$,
a raw coefficient is $m\in\Hom(M,\Phi N)$. It is identified with
an ordinary arrow by $m\mapsto\alpha_Nm$. The simultaneous operator on this
raw coefficient is
\[
 m\longmapsto\Phi(\alpha_N)\Phi(m)\alpha_M^{-1},
\]
not the ordinary-arrow formula~\eqref{mod:T}. Under the stated graph
identification it becomes the transported ordinary-arrow action.
The full multiobject identification, including permuted labels and
the last twisted bar face, is proved in
Lemma~\ref{tr:bar-contraction}; see
\eqref{tr:graph-identification}--\eqref{tr:graph-transport}.
Thus the paracyclic homotopy is transported through an actual
isomorphism of graph bimodules. In these coordinates $T$ commutes
with $b$, and multiplying $bH+Hb=1-T$ by $-T^{-1}$ proves
\eqref{mod:inverse-homotopy}. For permuted labels a common positive
power is used only to read their weights. The graph throughout is
canonically identified with the original $\Phi$-graph; neither a
power of $\Phi$ nor $\Phi^{-1}$ is substituted for it.
\end{proof}

\subsection{Compact normalization, including the frame labels}

For the measure calculation, modulus-one eigenvalues before twisting
are not enough: they must remain of modulus one after every compact
stabilizer twist. We establish a simultaneous normalization of the
deformation and frame data. Choose a faithful representation $R$ of
$\Ghat$ and include in the linear data the direct sum of its fibers at
all $d$ framed points. The $d$th power of the cyclic transporter acts
on these summands by the pure weight-zero closed-point Weil
transports. Hence the cyclic transporter itself has modulus-one
eigenvalues. Together with $V$ and $\mathfrak j$, these data give a
faithful finite-dimensional realization of the normalized algebraic
data group, a subgroup of a product of linear groups and
$K\rtimes\langle\text{cyclic permutation}\rangle$.

\begin{lemma}[Simultaneous compact normalization]
\label{mod:compact-normalization}
Let $J.e$ be a nilpotent orbit preserved by the normalized transporter
in Theorem~\ref{mod:quadratic-model}. After changing that transporter
by an element of $J$, choose a triple $(e,h,f)$ in this orbit and
write the normalized transporter as
\[
 u=t\,v_0.
\]
There are simultaneous compact forms of the group and frame data
such that:
\begin{enumerate}[label=\textup{(\roman*)}]
\item $t$ belongs to a compact group, preserves the compact form
      $J_{\mathrm{cpt}}$, and centralizes the triple;
\item $v_0$ is unipotent, centralizes $J$, and commutes with $t$;
\item for $L_e=Z_J(\mathrm{SL}_{2,e})$,
      $L_{e,\mathrm{cpt}}=L_e\cap J_{\mathrm{cpt}}$ is a maximal
      compact including its components, and every
      $k t v_0$, $k\in L_{e,\mathrm{cpt}}$, has modulus-one
      eigenvalues on all the normalized linear data.
\end{enumerate}
These assertions also hold on their tensor, symmetric, exterior,
subquotient, and triple-weight constructions. After the
orbit-fixing translation by $\lambda_e(\sqrt q)$, the semisimple
value of the frame transport around the $d$-cycle belongs to
\begin{equation}
\label{mod:cyclic-compact-value}
 \bigl\{[\lambda_e(q^{d/2})s]:
         s\in Z_{\Ghat}(\mathrm{SL}_{2,e})_{\mathrm{cpt}}\bigr\}
 =B_{\Ghat.e}.
\end{equation}
This holds for every $k\in L_{e,\mathrm{cpt}}$.
\end{lemma}

\begin{proof}
Write $u_1=s_1v_1$ for the Jordan decomposition of the original
normalized transporter in its algebraic normalizer. Both factors
preserve the data and normalize $J$. The semisimple part has compact
closure of its powers in the faithful realization just chosen.

We first remove the inner action of the unipotent part; this step
requires some care when $J$ is disconnected. For a complex reductive
group,
\begin{equation}
\label{mod:connected-automorphisms}
 \Aut(J)^\circ=\im\bigl(J^\circ\longrightarrow\Aut(J)\bigr).
\end{equation}
Indeed an infinitesimal automorphism is inner on the semisimple part
of $J^\circ$ and trivial on its central torus, whose character
lattice has discrete automorphism group. After subtracting this
inner action, its values on the other components give a cocycle of
the finite group $\pi_0(J)$ with values in
$\Lie Z(J^\circ)$. Averaging makes that cocycle a coboundary,
which is an inner infinitesimal action by the central torus. This
proves surjectivity of the differential of the inner-automorphism
map and hence~\eqref{mod:connected-automorphisms}.
The kernel of $J^\circ\to\Aut(J)$ is $Z(J)\cap J^\circ$, a diagonalizable central
group. Therefore the automorphism induced by $v_1$ has a unique
unipotent lift $a\in J^\circ$: take the unipotent part of any
lift; uniqueness follows because two such lifts differ by a central
semisimple element. Since $s_1$ commutes with $v_1$, uniqueness
gives $s_1as_1^{-1}=a$. The element $v_0=a^{-1}v_1$ is unipotent,
centralizes $J$, and commutes with $s_1$: here $v_1av_1^{-1}=a$
because the action of $v_1$ on $J$ is conjugation by $a$, so the
two unipotent factors defining $v_0$ commute. Changing the
transporter by $a^{-1}\in J$ thus leaves $s_1v_0$. Since $v_1$
acts on $\mathfrak j$ by an inner automorphism, it preserves every
$J$-orbit; the assumed stable orbit is therefore also preserved by
$s_1$.

Let $A_1$ be the Zariski closure of the group generated by $s_1$.
It is diagonalizable, normalizes $J$, and $M=JA_1$ is reductive:
its unipotent radical maps trivially to $M/J$ and would be a
connected normal unipotent subgroup of $J^\circ$. The
maximal-compact theorem for complex reductive groups, including
disconnected groups \cite[Definition~3.3 and
Theorems~3.6--3.7]{AdamsTaibi}, gives a
maximal compact $M_{\mathrm{cpt}}$ containing the compact closure
of $\langle s_1\rangle$. Its intersection with $J$ is a maximal
compact $J_{\mathrm{cpt}}$: conjugate a maximal compact of $J$
into $M_{\mathrm{cpt}}$ and use normality of $J$ and maximality
inside the compact intersection. The element $s_1$ preserves this
intersection. Averaging one
Hermitian form makes this whole compact group unitary on the
faithful simultaneous realization.

Choose the triple in $J.e$ so that its $\mathrm{SU}(2)$ lies in
$J_{\mathrm{cpt}}$. The two homomorphisms $\mathrm{SL}_2\to J$
given by this triple and its $s_1$-translate are $J$-conjugate,
because their nilpotents lie in the same orbit. Their restrictions
to $\mathrm{SU}(2)$ land in $J_{\mathrm{cpt}}$, so they are
already conjugate by an element $k_0\in J_{\mathrm{cpt}}$.
For completeness, if a complex conjugator is $b=kp$ in the polar
decomposition relative to $J_{\mathrm{cpt}}$
\cite[Lemma~3.5]{AdamsTaibi}, unitarity of the two
homomorphisms implies that $b^*b$ commutes with the first one. Its
positive square root $p$ therefore commutes with it, and $k$ is
the required compact conjugator. This works in every component.
Set $t=k_0s_1$, choosing the conjugator's direction so that $t$
centralizes the triple. Since $v_0$ centralizes $J$, it commutes
with $k_0$, with $t$, and with every compact twist from $L_e$.
The product $kt$ belongs to $M_{\mathrm{cpt}}$ and is unitary for
the same Hermitian form. Multiplication by the commuting unipotent
$v_0$ changes neither its eigenvalues nor its semisimple
character. The centralizer of $\mathrm{SU}(2)$ is stable under
the compact conjugation, so its compact intersection is a compact
real form of the full complex triple centralizer
\cite[Lemma~3.10]{AdamsTaibi}, including every component. This
proves \textup{(i)}--\textup{(iii)}. The same proof
applies to all the indicated functorial linear constructions.

It remains to check the cyclic frame value, rather than merely the
action on $V$. Embed the compact frame image of
$M_{\mathrm{cpt}}$ in a maximal compact of
$K\rtimes\langle\text{cyclic permutation}\rangle$. Its
intersection with $K$ is a product of compact forms of the $d$
factors. For $k\in L_{e,\mathrm{cpt}}$, the $i$th factor of
$(kt)^d$ is the product transport around the cycle starting at the
$i$th frame. It is compact and centralizes the triple in that
factor. Since $v_0$ commutes with $kt$, the corresponding factor
of $(ktv_0)^d$ has that compact element as its semisimple part.
The translation by $\lambda_e(\sqrt q)$ commutes with both
$kt$ and $v_0$, and $v_0$ commutes with $kt$. Thus, writing $w=kt$,
\[
 \bigl(\lambda_e(\sqrt q)wv_0\bigr)^d
   =\lambda_e(q^{d/2})w^dv_0^d.
\]
No commutation between $k$ and $t$ is needed. This proves
\eqref{mod:cyclic-compact-value}. Compact forms of the triple
centralizer are conjugate inside that centralizer, so the set of
ambient conjugacy values is independent of the compact form
chosen in this proof.
\end{proof}

Translation of a transporter by $J$ does not change an equivariant
calculation: it changes the chosen identification of the same
Frobenius functor. After the additional translation by
$\lambda_e(\sqrt q)$ the point $e$ is fixed, since singular
points originally scale by $q^{-1}$ and $e$ has triple weight $2$.
The resulting weights and the distinction between the original and
lifted gradings are recorded in Section~\ref{orb:section}.
